\section{Discussion}
\label{sec:discussion}

\begin{resume}
Le gain est modeste parce que le goulot n'\'etait pas l\`a o\`u on le croyait :
l'inf\'erence occupait d\'ej\`a 96 \`a 99\,\% du temps, donc parall\'eliser la
pr\'eparation CPU ne pouvait pas donner plus que quelques pour cent de
m\'ediane, plus un effet de remplissage. Les crit\`eres d'acceptation sont
franchis, mais la d\'ecision d'activer huit workers se paie d'une complexit\'e
r\'eelle, et plusieurs menaces p\`esent sur la g\'en\'eralisation des
r\'esultats. Ce chapitre les assume explicitement et en tire les le\c{c}ons.
\end{resume}

\subsection{Pourquoi pas un facteur huit}

La loi d'Amdahl \cite{amdahl} explique l'essentiel. Si une fraction $p$ du
temps est parall\'elisable, le gain maximal est $1/(1-p)$. Ici, la mesure des
phases donne $1-p \approx 0,97$ \`a $0,99$ : le gain maximal th\'eorique est
donc de 1 \`a 3\,\%. Le gain observ\'e, de 7 \`a 26\,\% selon la position,
d\'epasse ce plafond apparent, ce qui signifie que les vagues ne font pas que
parall\'eliser la collecte : elles am\'eliorent aussi la fa\c{c}on dont les
lots sont remplis. Le remplissage passe de 6,1 \`a 6,7 en finale, et c'est
cette hausse qui explique la plus grande partie du gain sur la position la
plus lente.

Il y a une seconde raison, plus profonde, pour laquelle le plafond est bas :
le GPU n'\'etait d\'ej\`a pas inoccup\'e pendant la collecte, pas plus que le
CPU ne l'\'etait pendant l'inf\'erence dans le r\'egime batch\'e. La boucle
mono-thread \emph{pr\'epare} le lot suivant pendant que le GPU \'evalue le
pr\'ec\'edent ? Non : dans l'architecture par vagues, les deux phases
alternent strictement, et c'est justement cette alternance que la production
continue supprimerait. Mais les mesures montrent que l'alternance co\^ute peu
(1,5 \`a 2,3\,\% d'attente de barri\`ere) parce que la collecte elle-m\^eme
est tr\`es courte face \`a l'inf\'erence. Autrement dit, l'architecture
batch\'ee \'etait d\'ej\`a proche de sa limite : le GPU est le goulot, et il
le reste.

\subsection{Ce qui a march\'e, ce qui a d\'e\c{c}u}

Ce qui a march\'e : la finale, o\`u le gain atteint $+22\,\%$ de m\'ediane et
o\`u le p95 de latence s'am\'eliore m\^eme de 12 \`a 17\,\%, parce que les
recherches lentes y sont justement celles qui b\'en\'eficiaient le plus du
remplissage complet. L'\'egalit\'e de qualit\'e tient sur 2\,500 puzzles, et
le self-play n'a pas boug\'e. La correction est d\'emontrable par invariants
et tests de mutation, pas seulement par empirisme.

Ce qui a d\'e\c{c}u : le milieu de jeu, o\`u le gain de m\'ediane reste sous
8\,\%, et le plafond d\`es deux workers, qui dit clairement que la
quatri\`eme et la huiti\`eme unit\'e de calcul ne servent presque \`a rien. La
tranche de 8 simulations, utilis\'ee par le protocole \`a temps \'egal, est
m\^eme un r\'egime o\`u huit workers reculent de 19\,\% : sur des arbres de
huit simulations, les collisions dominent. C'est une propri\'et\'e du couple
(taille de tranche, nombre de workers) qui m\'erite d'\^etre connue avant
d'activer le multic\oe{}ur dans un autre contexte que le bot.

\subsection{Menaces sur la validit\'e}

\begin{enumerate}[leftmargin=1.4em]
  \item \textbf{D\'erive d'horloge et thermique.} Les sessions varient de 15
  \`a 30\,\%. Tout a \'et\'e mesur\'e en alternance pour compenser, mais les
  valeurs absolues ne sont pas transf\'erables telles quelles d'une machine
  \`a l'autre.
  \item \textbf{Queue de latence sur 50 observations.} Le p95 y est estim\'e
  sur 50 valeurs par configuration ; c'est mieux que les 6 valeurs par
  passage de la campagne principale, mais cela reste un estimateur de queue
  imparfait. Le plan lui-m\^eme pr\'evient qu'un p95 instable exige de
  nouvelles observations.
  \item \textbf{\'Echantillon de puzzles.} Le banc est un sous-\'echantillon
  r\'egulier de 2\,500 puzzles sur 5\,000, stratifi\'e par rating. Il mesure
  le premier coup d'une position tactique, pas la force en partie compl\`ete.
  Le verdict de non-inf\'eriorit\'e est un argument de qualit\'e de recherche,
  pas une mesure Elo.
  \item \textbf{Arbre froid.} Toutes les mesures partent d'un arbre neuf. En
  partie r\'eelle, \code{update\_root} conserve le sous-arbre et une partie
  des simulations \'evite le r\'eseau. Par la loi d'Amdahl, le gain du
  multic\oe{}ur s'en trouverait r\'eduit. Le cas chaud n'est pas mesur\'e.
  \item \textbf{Table du bot.} Le bot utilise 4\,000\,000 d'entr\'ees (environ
  4,16\,Gio), les mesures 8\,192. Un taux de succ\`es de table plus \'elev\'e
  r\'eduirait encore la fraction acc\'el\'er\'ee.
  \item \textbf{R\'ep\'etition \'ecart\'ee.} La campagne de 2\,500 puzzles n'a
  pas \'et\'e r\'ep\'et\'ee ; l'incertitude r\'esiduelle de qualit\'e va
  jusqu'\`a $-0,8$ point. C'est la limite la plus importante du verdict, et
  elle est consign\'ee comme telle dans le rapport de r\'esultats.
  \item \textbf{Mesures GPU seulement.} Le multic\oe{}ur a \'et\'e mesur\'e
  sur le GPU. Sur CPU, le rapport entre collecte et inf\'erence est tr\`es
  diff\'erent, et le gain serait probablement plus important ; il n'a pas
  \'et\'e caract\'eris\'e.
\end{enumerate}

\subsection{Le co\^ut d'ing\'enierie}

Dix t\^aches, cinq composants, quatorze cibles de test et plus de deux cents
tests Python : le chantier est lourd pour un gain de l'ordre de 10\,\%. Faut-il
le regretter ? Trois arguments disent non. D'abord, la partie non
parall\'elisable \'etait d\'ej\`a optimis\'ee ; sans ce chantier, aucun gain
suppl\'ementaire n'\'etait disponible c\^ot\'e ordonnancement, et les mesures
le d\'emontrent. Ensuite, les composants produits (cache par snapshots,
\'evaluateur injectable, garde RAII, ex\'ecuteur persistant, collecteur
born\'e) sont exactement les fondations de la production continue : les
abstractions sont r\'eutilis\'ees, pas jet\'ees. Enfin, la discipline de test
a produit une base de r\'ef\'erence mesur\'ee (chronom\'etrages, harnais,
comparaison appari\'ee) qui sert d\'ej\`a \`a d'autres chantiers, comme
l'\'etude des co\^uts CPU-GPU.

Le risque, lui, est r\'eel : une structure concurrente qu'on ne comprend plus
est une dette. C'est pourquoi l'observabilit\'e a \'et\'e trait\'ee comme une
fonctionnalit\'e de premier ordre : compteurs de collisions, d'\'etats
\code{Pending}, de fuites de virtual loss, messages de parcours. Un arbre
corrompu se voit au lieu de se deviner.

\subsection{Le\c{c}ons}

\begin{enumerate}[leftmargin=1.4em]
  \item \textbf{Mesurer avant d'optimiser.} L'intuition initiale, « le CPU
  prépare pendant que le GPU calcule », \'etait fausse dans ce r\'egime : les
  deux alternaient, et le GPU dominait \`a 97\,\%. La chronom\'etrie par
  phase, ajout\'ee en premi\`ere t\^ache, a red\'efini l'objectif avant
  d'\'ecrire la moindre ligne de concurrence.
  \item \textbf{Tester la concurrence par la conception.} Sans
  ThreadSanitizer, la seule strat\'egie tenable est de r\'eduire la surface
  concurrente \`a quelques \'etats atomiques et une garde, puis de la
  harceler avec des hooks et des mutations.
  \item \textbf{Un faux r\'eseau vaut mille GPU.} Le \code{ControlledEvaluator}
  a rendu testables les erreurs, les blocages et les interleavings que le
  GPU ne reproduit jamais deux fois.
  \item \textbf{Comparer en alternance.} Toute comparaison de performance sur
  une machine \`a horloge variable exige un protocole appari\'e ; l'\'ecart
  entre l'ancien et le nouveau binaire \'etait invisible sans cela.
  \item \textbf{Attention aux d\'efinitions des quantiles.} Le p95 de
  latence est le p5 de d\'ebit ; s'\^etre tromp\'e de queue aurait fait
  refuser un r\'eglage qui passe les crit\`eres.
  \item \textbf{Un test peut \^etre faux.} Le premier test d'\'equivalence du
  pool chaud comparait la mauvaise r\'ef\'erence ; le refaire correctement a
  transform\'e un \'echec en preuve.
\end{enumerate}
